%------------------------------------------------------------------------------%
\begin{question}
  Determine uma equação vetorial para a reta
  \[
    \frac{x-1}{2}
    = \frac{y+3}{-1}
    = \frac{z-2}{4}
  \]
\end{question}

%------------------------------------------------------------------------------%
\begin{resolution}{Solução}
  Introduzimos um parâmetro comum
  \pause
  \[
    \frac{x-1}{2}
    = \frac{y+3}{-1}
    = \frac{z-2}{4} \pause
    \;=\; t
  \]
  \pause
  Assim,
  \[
    x =  1 + 2t \qquad\pause
    y = -3 -  t \qquad\pause
    z =  2 + 4t
  \]
  \pause
  Portanto
  \[
    (x, y, z) = (1, -3, 2) +t(2, -1, 4)
  \]

  \pause
  O ponto
  \(
    P = (1, -3, 2)
  \)
  pertence à reta

  \pause
  e
  \(
    \vec{v} = (2, -1, 4)
  \)
  é um vetor diretor
\end{resolution}

%------------------------------------------------------------------------------%
\endinput
